% ==============================================================================
% Lone Trail Theme Palette: theme_template
% Name: 配色规范范本 (Custom Theme Template & Guide)
% Designer: Rhine Lab Design System
% ==============================================================================
%
% 【如何自制新主题？】
% 1. 复制本文件，重命名为 theme_XXX.tex（例如 theme_143.tex）
% 2. 参考下方建议修改 RGB 纯数值（0~255）
% 3. 双击【换模板并编译】.bat，系统会自动识别并同步到主题列表！
%
% ==============================================================================
% 1. 核心主强调色（Primary Accent）
% ------------------------------------------------------------------------------
% 出现位置：封面主标题、大章节 Section 标题、页眉页脚装饰线条、标题盒边框
% 选色建议：必须有足够对比度（L* 30%~55%），切忌使用过浅或荧光色，保证浅底上清晰有力。
\definecolor{qiancheng}{RGB}{227, 101, 53}

% ==============================================================================
% 2. 标志性三色条带（The Lone Trail Tri-Stripe Motif）
% ------------------------------------------------------------------------------
% 出现位置：封面/正文三色分割线（separator_1）、封面左竖条、目录页边条
% 选色建议：三者明度与饱和度尽量均衡；lv（左冷/青绿）、huang（中过渡/橙黄）、hong（右点睛/砖红）。
\definecolor{lv}{RGB}{131, 187, 170}
\definecolor{huang}{RGB}{227, 164, 60}
\definecolor{hong}{RGB}{149, 44, 38}

% ==============================================================================
% 3. 纸张底色与正文排版（Canvas & Typography）
% ------------------------------------------------------------------------------
% qianbai（纸张主底色）：全书页面物理纸面底色。
% 选色建议：极高明度（RGB 推荐 242~252），可带微弱暖米、冷灰或纸浆白，切忌偏灰暗。
\definecolor{qianbai}{RGB}{248, 243, 231}

% mibai（暗纹与辅助底色）：背景几何网格点阵、水印暗纹。
% 选色建议：仅比 qianbai 深 3%~8%（微弱呼吸感即可），过深会喧宾夺主干扰正文阅读。
\definecolor{mibai}{RGB}{244, 237, 218}

% heise（正文字体颜色）：段落文本、公式、表格内容、参考文献。
% 选色建议：极低明度（RGB 推荐 0~35），推荐纯黑或带微冷调的炭黑，保证长时间阅读舒适。
\definecolor{heise}{RGB}{0, 0, 0}

% ==============================================================================
% 4. 六大数学定理环境盒子（Academic Box Ensembles）
% ------------------------------------------------------------------------------
% 出现位置：6 种 tcolorbox 数学盒子的标题栏底色与边框描边
% 选色建议：明度保持中等（L* 40%~60%），撑起白色标题字；六色色相尽量错开，
%           推荐冷暖相间，使读者扫一眼即可从颜色秒辨定理类型。
\definecolor{definition}{RGB}{245, 164, 0}   % 定义（推荐：橙黄 / 暖黄）
\definecolor{theorem}{RGB}{222, 4, 20}       % 定理（推荐：朱红 / 绯红，最醒目）
\definecolor{proposition}{RGB}{177, 0, 129}  % 命题（推荐：洋红 / 莓紫）
\definecolor{corollary}{RGB}{133, 194, 72}   % 推论（推荐：草绿 / 翠绿）
\definecolor{lemma}{RGB}{233, 93, 157}       % 引理（推荐：浅紫 / 樱粉）
\definecolor{example}{RGB}{0, 153, 216}      % 例题（推荐：天蓝 / 钴蓝）
